\subsection{Disclosure approach}
\label{sec:disc-approach}

The bug class spans many independently-maintained projects across
language ecosystems with no shared coordination forum. We propose
the following sequencing:

\begin{enumerate}
  \item \emph{Day 0:} private notification to each affected upstream
    maintainer via that project's security channel, with:
    a per-server STATE.md, the standardized reproducer
    (\fname{wave\{1,2\}/<ecosystem>/repro/}), the cProfile /
    pprof / async-profiler trace from the worst Mode B cell, and a
    cross-reference to this paper as the disclosure document.
  \item \emph{Day 60:} coordination call (asynchronous, via OSV / OSSF
    or per-maintainer email) to align on a fix-or-position
    decision.
  \item \emph{Day 90:} public disclosure of this paper, with
    upstream advisories filed for any project that has issued a
    fix or a security position by that date.
\end{enumerate}

The 90-day window is conservative for an amplification class with
no obvious in-the-wild exploitation (we have not seen evidence of
it being weaponized, though we have not actively searched
attacker telemetry).

\subsection{Predicted maintainer responses}
\label{sec:disc-predict}

We predict per-maintainer responses on the basis of public posture
and prior CVE handling:

\begin{description}
  \item[\fwid{hyper}] (Rust): declined. hyper issue
    \#4008~\cite{hyper_4008} is the explicit prior-art closure
    of this class as \fname{not\_planned}; we expect a polite
    re-affirmation pointing at
    \fname{http\_body\_util::Limited} and the proposed
    \fname{BodyExt::limit\_chunks(N)}.
  \item[\fwid{uvicorn} / \fwid{hypercorn} / \fwid{daphne}] (Python):
    likely receptive. Each is small in maintainer-bandwidth terms
    but has shown willingness to accept performance-related PRs.
    A batched-receive implementation could land in 2--3 months.
  \item[\fwid{Vert.x}] (JVM): likely receptive. The
    \fname{HttpObjectAggregator} primitive already exists; the
    decision is which API to expose so applications can opt into
    aggregation without breaking the streaming contract.
  \item[\fwid{Spring Boot} / Tomcat:] likely declined as
    server-side fix. Thread-per-connection plus the
    \fname{Servlet getInputStream} contract make per-recv batching
    invasive. Their security position will likely be "deploy
    behind nginx" plus tuning thread-pool sizing.
  \item[\fwid{node.js}]: unknown. The quadratic shape is a
    stronger argument than the linear one. The fix would have to
    land in \fname{lib/internal/streams/readable.js} and the
    Readable-stream behavior change has compatibility risk; we
    expect a long maintainer discussion rather than a quick patch.
  \item[\fwid{express} / \fwid{fastify}]: defer to Node.js.
  \item[\fwid{bandit}]: likely receptive. The quadratic is a
    clear bug in
    \fname{do\_read\_chunk!/4}; the fix is local and the project
    has accepted similar performance fixes.
  \item[\fwid{cowboy}]: likely receptive. The
    \fname{cowlib} maintainer accepted CVE-2026-7790; the
    per-chunk amplification is adjacent and the same author would
    handle it.
  \item[\fwid{nginx}, Apache, \fwid{HAProxy}]: unknown. These are
    C servers with substantial engineering teams; their reactions
    will depend on how the class is framed (security versus
    performance) and on how many of their downstream users push
    for the fix.
  \item[\fwid{Kestrel}]: nothing to disclose; already mitigated
    structurally. We notify Microsoft only as a courtesy to
    confirm the Pipelines design intent and to ask whether the
    pattern is documented as a security property.
\end{description}

\subsection{Coordination challenges}
\label{sec:disc-coord}

Cross-language coordinated disclosure has limited precedent. The
closest analogs:

\begin{itemize}
  \item Heartbleed (2014): single library (OpenSSL), wide
    downstream impact, coordinated via the OpenSSL team.
  \item Log4Shell (2021): single library (log4j), ecosystem-wide
    JVM impact, coordinated via Apache Software Foundation.
  \item The HTTP smuggling papers (Black Hat 2019, 2020):
    cross-vendor, coordinated via the researcher's pre-presentation
    private notification to each vendor.
\end{itemize}

The HTTP smuggling precedent is the closest fit: a single
researcher / small team notifies many vendors, each handles their
own response independently, and the public talk consolidates the
ecosystem picture. We follow that model.

\subsection{Funding and timeline}
\label{sec:disc-funding}

This work is self-funded; we have no budget for paid maintainer
coordination or for follow-up engineering with each project. We
expect some upstream maintainers to take longer than 90 days to
respond; in those cases we will publish their unanswered status
in a follow-up addendum.
